% ECONOMICS_PROBLEM: {"id": "O43-435", "title": "Long-run causal effects of institutions", "jel_code": "O43", "source_id": "OP-0130"}
% !TeX program = xelatex
\documentclass[11pt,a4paper]{article}
\usepackage[margin=23mm]{geometry}
\usepackage{fontspec}
\setmainfont{Latin Modern Roman}
\usepackage{amsmath,amssymb,mathtools}
\usepackage{xcolor,enumitem,booktabs,longtable,array,xurl,hyperref}
\definecolor{muted}{HTML}{53636C}
\hypersetup{colorlinks=true,urlcolor=blue,linkcolor=blue}
\setlength{\parindent}{0pt}
\setlength{\parskip}{5pt}
\newcommand{\R}{\mathbb R}
\newcommand{\E}{\mathbb E}
\newcommand{\N}{\mathbb N}
\newcommand{\ind}{\mathbf 1}
\newcommand{\Var}{\operatorname{Var}}
\newcommand{\Cov}{\operatorname{Cov}}
\newcommand{\supp}{\operatorname{supp}}
\DeclareMathOperator*{\argmin}{arg\,min}
\DeclareMathOperator*{\argmax}{arg\,max}
\newcommand{\entrymeta}[3]{{\sffamily\small\color{muted}\textbf{\texttt{#1}}\quad|\quad #2\par #3\par}}
\newcommand{\Problemref}[1]{\texttt{#1}}
\newcommand{\JELref}[2]{\texttt{#2} (JEL \texttt{#1})}
\begin{document}
\section*{Long-run causal effects of institutions}
\textbf{Problem ID:} \texttt{O43-435}\quad \textbf{JEL:} \texttt{O43}\par
% BEGIN ECONOMICS STATEMENT
\label{problem:OP-0130}
\label{atlas:prob:Q10}
{\sffamily\small\color{muted}\textbf{Q10} | Political economy | Editorial research problem family\par}
\subsection*{Definitions and assumptions}
Fix a target population, horizon $h$, institution vector $d\in\mathcal D$, and feasible reform $a$ with implementation and enforcement explicitly defined. Let $Y_h(a)$ be integrable development outcomes. Historical exposures $Z$ may affect institutions and other channels; a model class $\mathcal M$ specifies both pathways and selection into observed populations. The principal target is $\theta_h(a,a_0)=\mathbb E[Y_h(a)-Y_h(a_0)]$. An institutional index alone is not a unique intervention. An additional scalar linear sensitivity illustration uses centered $Y,D,Z$ with finite second moments, $\operatorname{Cov}(Z,D)\ne0$, and
\[
Y=\beta D+\gamma Z+u,\qquad \mathbb E[Zu]=0,\qquad |\gamma|\leq\Gamma.
\]
These restrictions define a separate homogeneous-effect model, not the unrestricted historical problem.
\subsection*{Formal research question}
Determine the identified set for the reform contrast over $m\in\mathcal M$ with the observed law, allowing declared direct historical effects and uncertainty in measurement. In the scalar illustration, report the sensitivity set
\[
\left\{\frac{\operatorname{Cov}(Z,Y)-\gamma\operatorname{Var}(Z)}{\operatorname{Cov}(Z,D)}:\gamma\in[-\Gamma,\Gamma]\right\}.
\]
Determine what additional designs and measurements narrow the reform set, and what transport restrictions are required to extrapolate local boundary or instrument estimates to a national, long-run reform.
\subsection*{Interpretation and scope}
The displayed sensitivity identity follows from the specified moment restriction and is not an open theorem. Its role is to expose dependence on exclusion. Without further assumptions, an instrument effect is not a global causal effect of all institutions. Historical slavery and forced labor can operate through culture, human capital and market access as well as institutions; their total effects do not identify a unique institutional mediation share. Geography and institutions can interact, so attribution need not be additive.
\subsection*{Source audit and status}
All three atlas citations are verified. Their IV and regression-discontinuity designs provide important setting-specific persistence evidence under maintained restrictions. The atlas's broad comparison cannot be read as an unresolved request to reproduce their estimates. The retained problem is an editorial identification and transportability agenda with explicit sensitivity to direct historical channels.
\subsection*{References}
{\footnotesize\raggedright
\begin{enumerate}[label={\textbf{[S\arabic*]}},leftmargin=3em]
\item Daron Acemoglu, Simon Johnson, and James A. Robinson (2001). \emph{The Colonial Origins of Comparative Development: An Empirical Investigation}. American Economic Review 91(5), 1369--1401.\par
\textit{Locator and source role:} Abstract and settler-mortality IV design; institutional interpretation depends on exclusion and measurement assumptions.\par
\url{https://www.aeaweb.org/articles?id=10.1257/aer.91.5.1369}
\item Nathan Nunn (2008). \emph{The Long-Term Effects of Africa's Slave Trades}. Quarterly Journal of Economics 123(1), 139--176.\par
\textit{Locator and source role:} Abstract and historical slave-export design; total historical effects need not be exclusively mediated by institutions.\par
\url{https://academic.oup.com/qje/article-abstract/123/1/139/1889789}
\item Melissa Dell (2010). \emph{The Persistent Effects of Peru's Mining Mita}. Econometrica 78(6), 1863--1903.\par
\textit{Locator and source role:} Abstract and historical-boundary regression discontinuity; local persistence evidence, not a universal institutional causal coefficient.\par
\url{https://onlinelibrary.wiley.com/doi/10.3982/ECTA8121}
\end{enumerate}
}

\subsection*{Common conventions: Source mathematical conventions}

Unless an entry states otherwise, agent and firm sets are finite. State and action spaces are measurable, strategies and outcome maps are measurable, probability laws are specified, and expectations exist for the targets under discussion. Infinite-horizon discount factors lie in $(0,1)$ and discounted objectives are finite. Norms, tolerances and complexity input sizes must be specified for computational comparisons. Constants or primitives that appear locally are inputs, not empirical facts asserted by this catalogue.

\textbf{Identification.} A statistical model $m\in\mathcal M$ implies an observable law $P_m$ and target $\theta(m)$. At law $P$, its identified set is
\[
\Theta(P;\mathcal M)=\{\theta(m):m\in\mathcal M,\ P_m=P\}.
\]
Point identification holds when this set is a singleton, or equivalently when $P_{m_1}=P_{m_2}$ implies $\theta(m_1)=\theta(m_2)$ for all compatible models. Sharp bounds mean bounds attained, or approached when the boundary is not attained, by compatible models. Writing this set is a definition, not a solution: the substantive task is to establish informative restrictions and derive or estimate the set. An unrestricted-confounding model may give only trivial bounds.

\textbf{Inference.} Identification, estimability and valid uncertainty quantification are separate questions. Fix $\alpha\in(0,1)$. A confidence procedure $C_n$ over a declared law class $\mathcal P$ has asymptotic uniform coverage at level $1-\alpha$ when
\[
\liminf_{n\to\infty}\inf_{P\in\mathcal P}\Pr_P\{\theta(P)\in C_n\}\geq1-\alpha.
\]
For set-identified targets the coverage event must be defined separately, for example coverage of the full identified set. If an entry uses identification-indexed models instead of uniquely indexed laws, the same coverage convention is applied over those models. Pointwise limits do not establish uniform validity, and asymptotic validity does not guarantee finite-sample precision. Sample sizes, dependence structures and weak-identification sequences are part of each application.

\textbf{Counterfactuals.} $Y_i(a)$ is a potential outcome under a specified intervention, including its timing, implementation and versions. It is not $Y_i$ conditional on voluntarily choosing $a$. A full intervention vector permits interference; shorthand scalar treatment notation does not silently impose no interference. Assignment, selection, attrition, anticipation and measurement must be specified. A claim about an instrument's local effect is not automatically a population policy effect.

\textbf{Economic equilibrium.} A counterfactual model includes best responses, constraints, feasibility, market clearing, state transitions and expectations consistent with the stated information. When several equilibria exist, either a declared selector is an input or the target is set-valued. Comparisons across policy regimes must use the stated selection convention. Reduced-form relationships are not presumed invariant after an equilibrium-changing intervention.

\textbf{Optimization and welfare.} Optimization is over the stated feasible set. If attainment is not established, $\sup$ or $\inf$ and $\varepsilon$-optimal choices replace an asserted maximizer or minimizer. For example, with value $V=\sup_{a\in\mathcal A}W(a)<\infty$, an $\varepsilon$-optimal set is $\{a\in\mathcal A:W(a)\geq V-\varepsilon\}$ for $\varepsilon>0$. Benefits, costs, risk and health or environmental outcomes can be aggregated only using declared units and conversion weights. Interpersonal utility weights, the treatment of fiscal transfers and foreign profits, discounting, and the behavioral welfare criterion are explicit inputs. A comparative-static derivative additionally requires existence and appropriate differentiability of the selected counterfactual.

\textbf{Sources.} Local labels [S1], [S2], and so forth restart in every question. Locators identify the relevant abstract, model, section or result at the level actually checked; a bibliographic or abstract-level verification is not represented as inspection of an entire proof. Working papers are distinguished from later publications and changing versions. Source summaries are paraphrased. All URL checks and editorial formulations refer to the audit date stated above.
% END ECONOMICS STATEMENT
\end{document}
